\section{Conclusions}
\label{sec:conclusions}

An analytical model has been developed for branching initiation from a
prescribed stationary V-shaped interfacial shear crack at the corner of a
V-shaped interface between two isotropic elastic materials. The two symmetric
crack branches have faces in frictionless contact, and the applied loading is
assumed insufficient to cause loss of local stability at their outer tips.
Local fracture near the corner is represented by a small tensile process zone
oriented along the bisector of the angular sector occupied by one of the
materials and modeled as a line of normal displacement discontinuity subjected
to a constant normal cohesive traction. Mellin transformation, asymptotic matching, and the Wiener--Hopf method yield
analytical expressions for the zone length $d_i$, its opening $\delta_i$,
and the zone-related energy parameter $J_i$.

Physical admissibility of this local configuration requires the simultaneous
conditions
\begin{equation*}
  Cg_2<0,\qquad CQ_i>0.
\end{equation*}
The first ensures compressive contact between the interfacial crack faces,
whereas the second requires a tensile normal field in the material in which a
process zone can form. For the baseline material pair
$E_1/E_2=0.5$, $\nu_1=\nu_2=0.3$, these conditions partition the angular
range into physically admissible branches with transitions at
$\alpha_1=12.9202^\circ$ and $\alpha_2=107.0897^\circ$, and uniquely
identify which of the two materials can contain a zone for a prescribed sign
of $C$. Calculations for other values of $E_1/E_2<1$ show a systematic shift
of these transition angles with elastic contrast.

The analytical relations establish a power-law dependence of the process-zone
size on the applied loading and the material resistance to separation. On each physically admissible branch, increasing the magnitude of the
normalized load $|\sigma'|$ increases $d_i/l$, $\delta_i'$, and $J_i'$; for fixed geometry, both $\delta_i$ and $J_i$ are proportional to
the zone length. Decreasing the cohesive traction $\sigma_i$ increases both
the zone length and its opening. On the broad physically admissible branches,
the angular dependencies of the three parameters have nearby but distinct
maxima because the proportionality factors relating $d_i/l$, $\delta_i'$,
and $J_i'$ themselves depend on angle. These maxima identify geometries at
which the largest local deformation and energy measures are attained for a
prescribed external load. Determining the angular dependence of the critical
load requires a separately specified critical material parameter and a
separate solution of the corresponding criterion at each angle.

Comparison of the intact corner, the configuration with the prescribed
interfacial shear crack, and the configuration containing both the crack and
the process zone shows how the order of the corner singularity changes. For the baseline material pair, within the physically admissible segments,
the configuration containing the interfacial shear crack has a substantially
stronger singularity than the intact corner, whereas adding the process zone
partially shields it; specifically,
\begin{equation*}
  |\lambda_0|<|\lambda_i|<|\lambda|.
\end{equation*}
Thus, the singularity itself creates conditions favorable to local fracture
but is not a sufficient branching criterion: physical admissibility, the
resistance to separation, and the applied loading must also be taken into
account.

Within the adopted local model, $J_i$ is interpreted as a zone-related
energy parameter. Given an independently specified critical material
parameter $J_{i\mathrm{c}}$, the condition $J_i=J_{i\mathrm{c}}$ determines
the critical applied load for initiation of opening-mode (Mode~I) branching.
Alternatively, the deformation criterion
$\delta_i=\delta_{i\mathrm{c}}$ may be used. For fixed geometry, the two
criteria describe the same one-parameter family of equilibrium states and are
equivalent only when the corresponding critical material parameters are
chosen consistently.

The quantitative results should be interpreted within the assumptions of the
model: $d_i\ll l$, plane strain, symmetric loading, frictionless contact, and
retention of the leading singular term of the outer field. Over part of the
angular range, the calculated ratio $d_i/l$ reaches approximately $0.18$;
near the corresponding maxima, the small-scale assumption therefore becomes
marginal and quantitative estimates should be treated with caution. The model addresses branching initiation only; the subsequent path and stable
propagation of the secondary crack lie outside the scope of the present
formulation.
